\begin{trapbox}

	\begin{description}[style=nextline, leftmargin=2.8em, labelsep=0.5em, itemsep=0.35em]
		\item[\textbf{\textcolor{CTrap}{易错点一（椭圆外条件）}}]
			认为椭圆内部点也可以引两条垂直切线。
		\item[\textbf{\textcolor{CTrap}{正确理解（切线存在前提）：}}]
			只有椭圆外的点才能作两条实切线，椭圆内点不能作切线。
		\item[\textbf{\textcolor{CTrap}{错因分析（忽略位置）：}}]
			忽略了切线存在要求点不在椭圆内部。
	\end{description}

	\medskip
	\begin{description}[style=nextline, leftmargin=2.8em, labelsep=0.5em, itemsep=0.35em]
		\item[\textbf{\textcolor{CTrap}{易错点二（参数混淆）}}]
			在蒙日圆公式中误写为 $x^2+y^2=a^2-b^2$ 或 $x^2+y^2=b^2-a^2$。
		\item[\textbf{\textcolor{CTrap}{正确理解（公式记忆）：}}]
			蒙日圆方程为 $x^2+y^2=a^2+b^2$，其中 $a,b$ 分别为椭圆长半轴和短半轴。
		\item[\textbf{\textcolor{CTrap}{错因分析（概念混淆）：}}]
			对 $a,b$ 的意义不清晰，或与其他结论混淆。
	\end{description}

	\medskip
	\begin{description}[style=nextline, leftmargin=2.8em, labelsep=0.5em, itemsep=0.35em]
		\item[\textbf{\textcolor{CTrap}{易错点三（遗漏特例）}}]
			只考虑切线斜率存在的情形，认为点 $P$ 不能在 $x$ 轴或 $y$ 轴上。
		\item[\textbf{\textcolor{CTrap}{正确理解（包含特例）：}}]
			当一条切线斜率不存在时，另一条斜率为 $0$，此时点 $P$ 仍在蒙日圆上，公式依然成立。
		\item[\textbf{\textcolor{CTrap}{错因分析（思维不全）：}}]
			未能对斜率不存在的情况进行验证。
	\end{description}

\end{trapbox}
